\documentclass[11pt]{article}
\usepackage[margin=1in]{geometry}
\usepackage{amsmath,amssymb,amsthm,mathtools}
\usepackage{hyperref}

\newtheorem{theorem}{Theorem}[section]
\newtheorem{lemma}[theorem]{Lemma}
\newtheorem{proposition}[theorem]{Proposition}
\newtheorem{corollary}[theorem]{Corollary}
\theoremstyle{definition}
\newtheorem{question}[theorem]{Question}
\newtheorem{convention}[theorem]{Convention}
\theoremstyle{remark}
\newtheorem{remark}[theorem]{Remark}

\newcommand{\Z}{\mathbb{Z}}
\newcommand{\Q}{\mathbb{Q}}
\newcommand{\F}{\mathbb{F}}
\newcommand{\hgtt}{\operatorname{ht}}
\newcommand{\Cop}{\mathcal{C}}
\newcommand{\End}{\operatorname{End}}
\newcommand{\rev}{\operatorname{rev}}
\newcommand{\cpl}{\operatorname{cpl}}

\title{Particle--hole conjugation is not Terwilliger's $\dagger$:\\
the symmetry that lifts is the one that buys nothing}
\author{Clio}
\date{17 September 2026}

\begin{document}
\maketitle

\begin{abstract}
Q151 asked whether particle--hole conjugation $\Cop$ of the fermionic Fock space is the
antiautomorphism $\dagger$ of Terwilliger \cite[\S4]{T2020}, in the hope that
\cite[\S4, \texttt{prop:threefourths}]{T2020} would transfer to the classification of
commuting ribbon operators and make half of the nonvanishing work free.

The answer is \textbf{no}, on two independent counts, and the reason is structural rather
than technical.  We prove (Theorem~\ref{thm:A}) that $\Cop$ conjugates the rank-$e$ shape
weight operator $R_e^{W}$ to $R_e^{W^\vee}$, where $W^\vee$ is $W$ precomposed with
reversal-and-complementation of the occupancy word; the proof is three lines on the bead
lattice and replaces a $1500$-ribbon computational verification.  But $X\mapsto \Cop
X\Cop$ is conjugation by an involution, hence an \emph{automorphism}, while Terwilliger's
$\dagger$ is by definition an \emph{antiautomorphism}, and it is characterised by
\emph{fixing} the two generators of a Leonard pair, which $\Cop$ does not do.  Separately,
$\texttt{prop:threefourths}$ turns out to say that any \emph{three of four} tridiagonality
conditions imply the fourth; it is not a statement about three quarters of anything, and
the resemblance to \texttt{cor:quarter} of \cite{Q150c2} is a collision of names.

What survives is sharper than what was asked for, and it is not what a first pass found.
The admissibility system of \cite[Rem.~\texttt{rem:sym}]{Q150c2} carries a group $(\Z/2)^2$
generated by reversal and complementation of words.  We show
(Theorems~\ref{thm:C},~\ref{thm:D}) that exactly the subgroup generated by the composite
$T=\cpl\rev$ is implemented by a Fock-space involution --- no invertible graded linear map
implements reversal alone (Proposition~\ref{prop:norev}) --- and that what this subgroup
buys splits sharply in two.  On the \emph{nowhere-zero} part of the commuting locus it acts
by a scalar $\kappa=(-1)^{e-1}\gamma(d/2)^{e/d}$ and buys nothing.  On the
\emph{single-shape} part --- the part that \cite[Lem.~8.2]{Q149} wrongly excluded and
\cite{Q150c1,Q150c2} recovered --- it acts as the free involution $w\mapsto w^{T}$ on ribbon
shapes and halves the classification, its fixed points being exactly the antipalindromes,
i.e.\ exactly the strong-condition family of
\cite[Thm.~\texttt{thm:singleton2}]{Q150c2}.  So the Terwilliger \emph{mechanism} does
transfer; it is just not delivered by a $\dagger$, and it delivers only where the older
classification was blind.
\end{abstract}

\tableofcontents

\section{Setting and statement}

We use the normalisation of \cite{Q149} throughout.

\begin{convention}[Maya sets, moves, occupancy words]\label{conv:setup}
For a partition $\lambda$ put $M(\lambda)=\{\lambda_j-j:j\ge1\}\subset\Z$, a set containing
all sufficiently negative integers and bounded above.  A \emph{legal rank-$e$ move} in a
Maya set $M$ is a pair $b\mapsto b+e$ with $b\in M$ and $b+e\notin M$; it corresponds to
adding a connected $e$-ribbon to $\lambda$.  Its \emph{occupancy word} is
\[
   u_M(b,e)=(u_1,\dots,u_{e-1})\in\{0,1\}^{e-1},\qquad u_i=[\,b+i\in M\,].
\]
A \emph{shape weight} of rank $e$ over a commutative ring $K$ is a function
$W:\{0,1\}^{e-1}\to K$; it defines
\[
   R_e^{W}\,|M\rangle=\sum_{\substack{b\in M\\ b+e\notin M}} W\bigl(u_M(b,e)\bigr)\,
   \bigl|(M\setminus\{b\})\cup\{b+e\}\bigr\rangle,
\]
equivalently $R_e^{W}s_\lambda=\sum_\mu W(u)\,s_\mu$ over $e$-ribbons $\mu/\lambda$.
\end{convention}

\begin{convention}[the three word involutions]\label{conv:invol}
On $\{0,1\}^{n}$ define
\[
  (\rev u)_i=u_{n+1-i},\qquad (\cpl u)_i=1-u_i,\qquad
  u^{T}:=\cpl\rev\,u=\rev\cpl\,u,\quad (u^{T})_i=1-u_{n+1-i}.
\]
For a shape weight $W$ of rank $e$ we write $W^\vee(u):=W(u^{T})$, the weight precomposed
with reversal-and-complementation on $\{0,1\}^{e-1}$.  By
\cite[Lem.~\texttt{lem:transpose}]{Q149}, $u\mapsto u^{T}$ is transposition of the ribbon.
\end{convention}

\begin{convention}[particle--hole conjugation]\label{conv:C}
For a Maya set $M$ put
\[
   \Cop(M)\;:=\;\{\,-1-x \;:\; x\in\Z\setminus M\,\}.
\]
This is again a Maya set of charge $0$, and $\Cop^2=\mathrm{id}$.  We also write $\Cop$ for
the induced linear involution $|M\rangle\mapsto|\Cop M\rangle$ of the charge-$0$ Fock
space.
\end{convention}

\begin{lemma}\label{lem:Ctrans}
$\Cop\bigl(M(\lambda)\bigr)=M(\lambda')$: particle--hole conjugation is transposition of
partitions.
\end{lemma}

\begin{proof}
Classical.  Reading $\epsilon_x=[x\in M(\lambda)]$ left to right along $\Z$ gives the
boundary lattice path of $\lambda$, with $\epsilon_x=1$ a vertical step and $\epsilon_x=0$
a horizontal step, normalised so that charge $0$ puts the corner of the staircase between
$-1$ and $0$.  Transposing $\lambda$ reflects the diagram in the main diagonal, which
reverses the path about that point and exchanges the two step types; on $\epsilon$ this is
exactly $\epsilon_x\mapsto 1-\epsilon_{-1-x}$, i.e.\ $M\mapsto\Cop(M)$.
Verified independently on all $97$ partitions of size $\le9$ (Check~V1 of \S\ref{sec:ver}).
\end{proof}

\section{Theorem A: what $\Cop$ does to a ribbon operator}

\begin{theorem}\label{thm:A}
Let $e\ge1$ and let $W$ be a shape weight of rank $e$ over any commutative ring.  Then
\[
   \Cop\,R_e^{W}\,\Cop \;=\; R_e^{W^\vee},\qquad W^\vee(u)=W(u^{T}).
\]
\end{theorem}

\begin{proof}
Fix a Maya set $M$ and write $M^{c}=\Cop(M)$, so that for every $x\in\Z$
\begin{equation}\label{eq:flip}
   x\in M \iff -1-x\notin M^{c}.
\end{equation}

\emph{Moves correspond.}  Let $b\mapsto b+e$ be a legal move in $M$ and put
$b^{c}:=-1-b-e$.  By \eqref{eq:flip}, $b\in M$ gives $-1-b\notin M^{c}$, i.e.\
$b^{c}+e\notin M^{c}$; and $b+e\notin M$ gives $-1-b-e\in M^{c}$, i.e.\ $b^{c}\in M^{c}$.
So $b^{c}\mapsto b^{c}+e$ is a legal rank-$e$ move in $M^{c}$.  The assignment
$b\mapsto -1-b-e$ is an injection, and applying the same argument to $M^{c}$ (using
$\Cop^{2}=\mathrm{id}$) produces the inverse; so it is a bijection between the legal
rank-$e$ moves of $M$ and those of $M^{c}$.

\emph{Words reverse-and-complement.}  For $1\le i\le e-1$,
\[
   u_{M^{c}}(b^{c},e)_i=[\,b^{c}+i\in M^{c}\,]=[\,-1-(b+e-i)\in M^{c}\,]
   \overset{\eqref{eq:flip}}{=}[\,b+e-i\notin M\,]=1-u_M(b,e)_{e-i},
\]
where $e-i$ runs over $1,\dots,e-1$ as $i$ does.  That is exactly
$u_{M^{c}}(b^{c},e)=u_M(b,e)^{T}$.

\emph{Targets correspond.}  Write $M'=(M\setminus\{b\})\cup\{b+e\}$.  The set $M'$ differs
from $M$ only at $b$ and $b+e$, with both memberships flipped; by \eqref{eq:flip}, $\Cop
M'$ differs from $M^{c}$ only at $-1-b=b^{c}+e$ and $-1-b-e=b^{c}$, again with both
flipped.  Hence $\Cop M'=(M^{c}\setminus\{b^{c}\})\cup\{b^{c}+e\}$, the target of the
corresponding move.

Now compute.  Since $\Cop^{2}=\mathrm{id}$,
\[
  \Cop R_e^{W}\Cop\,|M\rangle=\Cop R_e^{W}|M^{c}\rangle
   =\Cop\!\!\sum_{b^{c}}\! W\bigl(u_{M^{c}}(b^{c},e)\bigr)\bigl|\cdots\bigr\rangle
   =\sum_{b} W\bigl(u_M(b,e)^{T}\bigr)\,\bigl|M'\bigr\rangle
   =R_e^{W^\vee}|M\rangle. \qedhere
\]
\end{proof}

\begin{remark}[a computational verification becomes a proof]
\cite[Lem.~\texttt{lem:transpose}]{Q149} --- transposition acts on occupancy words by
reversal-and-complementation --- was recorded there as \emph{verified} against $1500$
ribbons, with no derivation.  The middle paragraph of the proof above is a derivation, and
it never mentions Young diagrams: \eqref{eq:flip} does all the work.  Combined with
Lemma~\ref{lem:Ctrans} it upgrades \texttt{lem:transpose} from \texttt{computed} to
\texttt{proved}.
\end{remark}

\begin{corollary}\label{cor:aut}
$X\mapsto \Cop X\Cop$ is an algebra \emph{automorphism} of $\End$ of the Fock space, and it
carries $R_e^{W}\mapsto R_e^{W^\vee}$ for every $e$ and $W$.  In particular, for weights of
ranks $e,f$,
\[
   \Cop\,\bigl(R_e^{W}R_f^{\bar W}\bigr)\,\Cop \;=\; R_e^{W^\vee}\,R_f^{\bar W^\vee},
\]
with the factors in the \emph{same} order.
\end{corollary}

\begin{proof}
Conjugation by an involution is multiplicative: $\Cop XY\Cop=(\Cop X\Cop)(\Cop Y\Cop)$.
Apply Theorem~\ref{thm:A} to each factor.
\end{proof}

\section{Q151, answered}

\subsection{Terwilliger's $\dagger$, from its definition}

We quote \cite[\S4]{T2020} verbatim in substance.  Let $\Phi=(A;\{E_i\};A^*;\{E^*_i\})$ be
a pre Leonard system on $V$ with $\dim V=d+1<\infty$, and assume
$E^*_iAE^*_j=0$ for $|i-j|>1$ and $\ne0$ for $|i-j|=1$.  An \emph{antiautomorphism} of
$\End(V)$ is an $\F$-linear bijection $\zeta$ with $(XY)^\zeta=Y^\zeta X^\zeta$ for all
$X,Y$.  Then \cite[Lem.~\texttt{lem:antiaut}]{T2020}:
\begin{quote}
there exists a \emph{unique antiautomorphism} $\dagger$ of $\End(V)$ that \emph{fixes each
of $A,A^*$}; it fixes everything in $\mathcal D$ and $\mathcal D^*$, fixes each $E_i,E^*_i$,
and $(X^\dagger)^\dagger=X$.
\end{quote}

Two features of this definition are load-bearing, and $\Cop$ has neither.

\begin{proposition}[variance]\label{prop:variance}
$X\mapsto\Cop X\Cop$ is multiplicative, not antimultiplicative, on the algebra generated by
the ribbon operators.  Concretely, for generic weights $W,\bar W$ of ranks $e,f$,
\[
   \Cop\bigl(R_e^{W}R_f^{\bar W}\bigr)\Cop \;=\; R_e^{W^\vee}R_f^{\bar W^\vee}
   \;\ne\; R_f^{\bar W^\vee}R_e^{W^\vee},
\]
the two candidates differing by $[\,R_e^{W^\vee},R_f^{\bar W^\vee}\,]$.
\end{proposition}

\begin{proof}
The equality is Corollary~\ref{cor:aut}.  For the inequality it suffices to exhibit one
pair of ranks and one $\lambda$ with $[R_e^{W^\vee},R_f^{\bar W^\vee}]s_\lambda\ne0$ for
generic $W,\bar W$; by \cite[Thm.~\texttt{thm:main}]{Q149} the commuting locus is a proper
closed subvariety of weight space whenever it is nonempty, so generic weights do not
commute.  Check~V3 of \S\ref{sec:ver} exhibits the witnesses explicitly.
\end{proof}

\begin{remark}[naming the separator in advance]\label{rem:sep}
An automorphism and an antiautomorphism agree on every product of length $\le1$.  No
computation with a \emph{single} ribbon operator can tell them apart, and --- the point
that had to be written down before running anything --- \emph{on the commuting locus the
two candidates coincide}, because $[A,B]=0$ makes $AB=BA$.  A test performed with
Murnaghan--Nakayama weights, or with any pair drawn from the locus of
\cite[Thm.~\texttt{thm:main}]{Q149}, is constant in the direction it is supposed to
measure.  The separating coordinate is a product of two operators with
\emph{independent generic} weights.  Check~V3 runs there and reports how many $\lambda$
actually separate, so that the control is not silently empty.
\end{remark}

\begin{proposition}[the fixing property]\label{prop:fixing}
$\dagger$ is characterised among antiautomorphisms by fixing each of $A,A^*$.  The
analogous statement fails for $\Cop$: by Theorem~\ref{thm:A}, $\Cop R_e^{W}\Cop=R_e^{W}$ if
and only if $W^\vee=W$, i.e.\ $W$ is invariant under ribbon transposition, and for $e\ge2$
that is a \emph{proper} subspace of weight space, of dimension $2^{e-2}$ for $e$ even and
$2^{e-2}+2^{(e-3)/2}$ for $e$ odd.
\end{proposition}

\begin{proof}
Put $n=e-1$.  The fixed space of $W\mapsto W\circ T$ on functions
$\{0,1\}^{n}\to K$ has dimension equal to the number of orbits of $T$ on $\{0,1\}^{n}$,
namely $\tfrac12\bigl(2^{n}+\#\mathrm{Fix}(T)\bigr)$.  A word is fixed iff
$u_i+u_{n+1-i}=1$ for all $i$.  If $n$ is odd the index $i=(n+1)/2$ gives $2u_i=1$, which
is impossible, so $\#\mathrm{Fix}(T)=0$; if $n$ is even, $u$ is free on
$i\le n/2$ and forced on the rest, so $\#\mathrm{Fix}(T)=2^{n/2}$.  Hence
\[
   \dim\{W:W^\vee=W\}=\begin{cases}2^{e-2}, & e \text{ even},\\[2pt]
   2^{e-2}+2^{(e-3)/2}, & e \text{ odd},\end{cases}
\]
which is $<2^{e-1}$ for every $e\ge2$.
\end{proof}

\subsection{\texttt{prop:threefourths} is a collision of names}

\begin{remark}\label{rem:collision}
\cite[\texttt{prop:threefourths}]{T2020} reads, in full: given four conditions (i)--(iv)
asserting one-sided tridiagonality of $A$ with respect to $\{E^*_i\}$ and of $A^*$ with
respect to $\{E_i\}$, \emph{if at least three of \textrm{(i)--(iv)} hold then all four
hold}, i.e.\ the pre Leonard system is a Leonard system.  ``Three fourths'' counts
\emph{conditions}, not shapes.  \texttt{cor:quarter} of \cite{Q150c2} says that a nonzero
commuting weight is nowhere zero or vanishes on more than three quarters of the rank-$d$
shapes.  The two statements share a fraction and nothing else.  There is no theorem to
transfer.
\end{remark}

\begin{remark}[what the mechanism actually is, and why it still does not transfer]
\label{rem:mechanism}
The \emph{proof} of \texttt{prop:threefourths} is one line and is the thing worth naming:
$\dagger$ is an antiautomorphism fixing $A,A^*$, hence $(E_iA^*E_j)^\dagger=E_jA^*E_i$,
hence condition (iii) holds if and only if (iv) does.  The pattern is: \emph{an involution
that pairs up the nonvanishing conditions makes half of them free.}  That is exactly the
hope Q151 was built on, and Theorem~\ref{thm:C} below is the reason it fails here: the
involution available to us acts on the commuting locus by a scalar, so the pairing it
induces on nonvanishing conditions is the trivial pairing, and half of nothing is nothing.
\end{remark}

\begin{theorem}[Q151]\label{thm:Q151}
Particle--hole conjugation $\Cop$ is not Terwilliger's $\dagger$.  It fails
antimultiplicativity (Proposition~\ref{prop:variance}) and it fails the fixing property
that characterises $\dagger$ (Proposition~\ref{prop:fixing}); and the proposition whose
transfer motivated the question is about a different subject
(Remark~\ref{rem:collision}).
\end{theorem}

\section{Theorem C: which symmetry lifts, and what it is worth}

\begin{remark}[the locus I must work over]\label{rem:which}
\cite[Thm.~\texttt{thm:main}]{Q149} describes the commuting locus as the multiplicative
$\gcd$-periodic characters $W^{\gamma}$.  That description is \emph{superseded}:
\cite[Cor.~\texttt{cor:main0}]{Q150c1} shows that \cite[Lem.~8.2]{Q149} --- nonvanishing of a
commuting weight --- is false for $\gcd(e,f)\ge3$, and
\cite[Thm.~\texttt{thm:main0}]{Q150c1} replaces the classification by: for
$W\not\equiv0\not\equiv\bar W$ and $d=\gcd(e,f)$,
\begin{equation}\label{eq:locus}
   [R_e^{W},R_f^{\bar W}]=0\iff
   \sigma=c\,\tau^{\otimes(e/d)},\quad \bar\sigma=\bar c\,\tau^{\otimes(f/d)},\quad
   \sigma\ \text{satisfies}\ (2\mathrm{B}^{*}_e),
\end{equation}
for some nonzero rank-$d$ shape weight $\tau$ over $F$, where $\sigma(u)=(-1)^{|u|}W(u)$ and
$\otimes$ is concatenation of shape weights \cite[Conv.~\texttt{conv:otimes}]{Q150c1}.  The
$W^{\gamma}$ are exactly the nowhere-zero points of \eqref{eq:locus}; the points with zeros
are the new ones, and they are where the answer to Q151 stops being trivial.  A first pass at
this section, written against \cite[Thm.~\texttt{thm:main}]{Q149}, concluded the opposite of
Theorem~\ref{thm:D}(b) below.
\end{remark}

\cite[Rem.~\texttt{rem:sym}]{Q150c2} records that admissibility of a support system is
preserved by $\cpl$ and by $\rev$, and that the composite is ribbon transposition $T$; the
remark after \cite[Cor.~\texttt{cor:48}]{Q150c2} records that the list of admissible single
shapes is stable under $T$.  Both are statements about \emph{words}.  Theorem~\ref{thm:A}
promotes one of them to a statement about \emph{operators}, and that is what makes the
accounting below possible.

\begin{lemma}[the involutions pass through $\otimes$]\label{lem:otimes}
Let $\tau$ be a shape weight of rank $d$, $k\ge1$, and $g\in\{\rev,\cpl,T\}$ acting on words
of the appropriate length.  Then
$\bigl(\tau^{\otimes k}\bigr)\circ g=\bigl(\tau\circ g\bigr)^{\otimes k}$.
\end{lemma}

\begin{proof}
Write $n=kd-1$.  In $(\tau^{\otimes k})(y)=\prod_{i=1}^{k}\tau\bigl(y_{[(i-1)d+1,\,id-1]}\bigr)$
the $k-1$ letters in positions $d,2d,\dots,(k-1)d$ are ignored.  Under
$i\mapsto n+1-i=kd-i$, position $(i-1)d+j$ with $1\le j\le d-1$ goes to $(k-i)d+(d-j)$, i.e.\
block $i$ goes to block $k+1-i$ with its letters reversed, and the ignored positions $id$ go to
the ignored positions $(k-i)d$.  So $\rev$ permutes the blocks and reverses each; since the
product over blocks is a product of scalars it is insensitive to the order of the blocks,
giving $(\tau^{\otimes k})\circ\rev=(\tau\circ\rev)^{\otimes k}$.  For $\cpl$ the block
positions are fixed and each block is complemented.  Compose for $T$.
Verified symbolically for all $d\le4$, $k\le3$ with $kd-1\le9$ ($1368$ identities, $0$
failures; Check~V7).
\end{proof}

\begin{lemma}[the two-bead condition is $G$-stable]\label{lem:2Bstable}
Let $\sigma$ be a shape weight of rank $m$ satisfying $(2\mathrm{B}^{*}_m)$.  Then
$\sigma\circ\rev$ and $\sigma\circ\cpl$ satisfy it, hence so does $\sigma\circ T$.
\end{lemma}

\begin{proof}
Write the condition as: for all $\delta\in[1,m-1]$ and all $\pi,\varsigma\in\{0,1\}^{\delta-1}$,
$\pi'\in\{0,1\}^{m-\delta-1}$,
$\sigma(\pi0\pi')\sigma(\pi'0\varsigma)=\sigma(\pi1\pi')\sigma(\pi'1\varsigma)$.

\emph{Complementation.}  $(\sigma\circ\cpl)(\pi b\pi')=\sigma(\bar\pi\,\bar b\,\bar\pi')$, so
the $b$-side for $\sigma\circ\cpl$ at $(\delta;\pi,\pi',\varsigma)$ equals the $\bar b$-side for
$\sigma$ at $(\delta;\bar\pi,\bar\pi',\bar\varsigma)$.  The hypothesis equates the two sides for
$\sigma$, hence for $\sigma\circ\cpl$.

\emph{Reversal.}  $(\sigma\circ\rev)(\pi b\pi')=\sigma(\tilde\pi'\,b\,\tilde\pi)$ where
$\tilde{\ }$ is reversal.  Hence
\[
  (\sigma\circ\rev)(\pi b\pi')\,(\sigma\circ\rev)(\pi' b\varsigma)
  =\sigma(\tilde\pi'\,b\,\tilde\pi)\,\sigma(\tilde\varsigma\,b\,\tilde\pi')
  =\sigma\bigl(\tilde\varsigma\,b\,\tilde\pi'\bigr)\,\sigma\bigl(\tilde\pi'\,b\,\tilde\pi\bigr),
\]
the last step being commutativity of the ring.  The right-hand side is the $b$-side of the
condition for $\sigma$ at the \emph{same} $\delta$, with
$(\pi,\pi',\varsigma):=(\tilde\varsigma,\tilde\pi',\tilde\pi)$ --- the lengths match, since
$|\tilde\varsigma|=|\tilde\pi|=\delta-1$ and $|\tilde\pi'|=m-\delta-1$.  Apply the hypothesis.
Verified exhaustively over $\mathbb F_5$ for $m=3,4$ (all $25$ resp.\ $81$ solutions; $0$
failures; Check~V8).
\end{proof}

\begin{theorem}[the symmetry group acts, and $T$ is the part an operator sees]\label{thm:C}
Let $1\le e<f$, $d=\gcd(e,f)$.  The group $G=\langle\rev,\cpl\rangle\cong(\Z/2)^2$ acts on the
commuting locus \eqref{eq:locus} by $\tau\mapsto\tau\circ g$ on the rank-$d$ seed.  Under this
action the subgroup $\langle T\rangle$ is implemented on operators: by Theorem~\ref{thm:A},
\[
   \Cop\,R_e^{W}\,\Cop=R_e^{W^{\vee}},\qquad \Cop\,R_f^{\bar W}\,\Cop=R_f^{\bar W^{\vee}},
\]
and $(W,\bar W)$ commute if and only if $(W^{\vee},\bar W^{\vee})$ do.
\end{theorem}

\begin{proof}
Lemma~\ref{lem:otimes} shows $g$ carries $\tau^{\otimes k}$ to $(\tau\circ g)^{\otimes k}$, so
the shape $\sigma=c\,\tau^{\otimes(e/d)}$, $\bar\sigma=\bar c\,\tau^{\otimes(f/d)}$ of
\eqref{eq:locus} is preserved with seed $\tau\circ g$, up to the global sign $(-1)^{m-1}$
incurred by $\cpl$ in passing between $W$ and $\sigma$ (which is a unit and does not affect
\eqref{eq:locus}).  Lemma~\ref{lem:2Bstable} preserves $(2\mathrm B^{*}_e)$.  The operator
statement is Corollary~\ref{cor:aut} applied to both factors.
\end{proof}

\begin{theorem}[the two regimes]\label{thm:D}
\emph{(a) The nowhere-zero part.}  Let $\tau$ be nowhere zero, so that
$W=\alpha W^{\gamma}_e$, $\bar W=\beta W^{\gamma}_f$ in the notation of
\cite[Conv.~\texttt{conv:gamma}]{Q149}.  Then
\[
  W^{\gamma}_e\circ\rev=W^{\gamma^{-1}}_e,\qquad
  W^{\gamma}_e\circ T=\kappa\,W^{\gamma}_e,\quad
  \kappa=(-1)^{e-1}\Gamma,\quad
  \Gamma=\begin{cases}1,& d \text{ odd},\\ \gamma(d/2)^{e/d}\in\{\pm1\},& d\text{ even}.
  \end{cases}
\]
So $T$ acts by a \emph{scalar} --- trivially, modulo the scalars $\alpha,\beta$ that
\eqref{eq:locus} quotients by --- while $\rev$ acts by $\gamma\mapsto\gamma^{-1}$, which moves
points if and only if $d\ge3$.

\emph{(b) The single-shape part.}  Let $\tau=c\,\mathbb 1_w$ with $w\in\{0,1\}^{d-1}$
satisfying \textnormal{(NB)} of \cite[Def.~\texttt{def:border}]{Q150c2} --- that is, writing
$n=d-1$ and calling $p\in[0,n-1]$ a \emph{border length} when $w_{[1,p]}=w_{[n-p+1,n]}$
(so $p=0$ always is one), $w_{p+1}\ne w_{n-p}$ for every border length $p$.  Then
$\tau\circ T=c\,\mathbb 1_{w^{T}}$, and $w^{T}$ again satisfies \textnormal{(NB)}.  Here $T$
does \emph{not} act by a scalar: the smallest instance is $d=4$, where
$w=110\mapsto w^{T}=100$, i.e.\ the ribbon shape $(2,1,1)$ is exchanged with $(3,1)$.  The
fixed points of $T$ are exactly the \emph{antipalindromes}, hence exactly the $w$ for which
$\mathbb 1_w$ satisfies the strong condition
$(2\mathrm B^{**}_d)$ \cite[Thm.~\texttt{thm:singleton2}]{Q150c2}.
\end{theorem}

\begin{proof}
(a) Write $\sigma^{\gamma}_e(u)=\prod_j\gamma(j\bmod d)^{u_j}$, so
$W^{\gamma}_e=(-1)^{|u|}\sigma^{\gamma}_e$.  Then
$\sigma^{\gamma}_e(\rev u)=\prod_{j}\gamma(j)^{u_{e-j}}=\prod_i\gamma(e-i)^{u_i}$, and $d\mid e$
gives $\gamma(e-i\bmod d)=\gamma(-i\bmod d)=\gamma(i)^{-1}$, so this is
$\sigma^{\gamma^{-1}}_e(u)$; as $|\rev u|=|u|$, the first identity follows.  Also
$\sigma^{\gamma}_e(\cpl u)=\prod_j\gamma(j)^{1-u_j}=\Gamma\,\sigma^{\gamma^{-1}}_e(u)$ and
$|\cpl u|=e-1-|u|$, so $W^{\gamma}_e\circ\cpl=(-1)^{e-1}\Gamma\,W^{\gamma^{-1}}_e$; composing
gives $W^{\gamma}_e\circ T=(-1)^{e-1}\Gamma\,W^{\gamma}_e$.  For $\Gamma$: grouping
$j=1,\dots,e-1$ into $e/d$ full residue systems and using $\gamma(0)=1$,
$\Gamma=\bigl(\prod_{\delta\in\Z/d}\gamma(\delta)\bigr)^{e/d}$; the relation
$\gamma(\delta)\gamma(-\delta)=1$ kills every two-element orbit, leaving $\gamma(0)=1$ and,
for $d$ even, $\gamma(d/2)$, which squares to $1$.  Finally $\gamma\mapsto\gamma^{-1}$ is the
identity on the whole parameter group iff that group is $2$-torsion; by
\cite[Cor.~\texttt{cor:size}]{Q149} it is
$(F^{\times})^{\lfloor(d-1)/2\rfloor}\times\mu_2(F)^{[2\mid d]}$, which is $2$-torsion iff
$\lfloor(d-1)/2\rfloor=0$, i.e.\ $d\le2$.

(b) $(\mathbb 1_w\circ T)(u)=[\,u^{T}=w\,]=[\,u=w^{T}\,]$ since $T^{2}=\mathrm{id}$.  That
$w^{T}$ satisfies (NB) is Lemma~\ref{lem:2Bstable} together with
\cite[Thm.~\texttt{thm:singleton}]{Q150c2}; directly, $T$ preserves border lengths
($w^{T}_{[1,p]}$ is the complemented reverse of $w_{[n-p+1,n]}$ and vice versa) and
$w^{T}_{p+1}=1-w_{n-p}$, $w^{T}_{n-p}=1-w_{p+1}$, so the clause $w_{p+1}\ne w_{n-p}$ is
carried to itself.  The witness $d=4$ is Check~V9.  For the fixed points:
$w=w^{T}$ says $w_i=1-w_{n+1-i}$ for all $i$, i.e.\ $w_i\ne w_{n+1-i}$, which is the
antipalindrome condition of \cite[Thm.~\texttt{thm:singleton2}]{Q150c2}.
\end{proof}

\begin{corollary}[the parity, explained]\label{cor:parity}
$T$ has no fixed point on $\{0,1\}^{n}$ for $n$ odd, and exactly $2^{n/2}$ for $n$ even.
Hence the $(2\mathrm B^{**})$ single-shape family is empty for $d$ even and has
$2^{(d-1)/2}$ members for $d$ odd --- recovering the count of
\cite[Thm.~\texttt{thm:singleton2}]{Q150c2} from the fixed-point count of particle--hole
conjugation alone.
\end{corollary}

\begin{proof}
$w=w^{T}$ reads $w_i+w_{n+1-i}=1$.  For $n$ odd the index $i=(n+1)/2$ demands $2w_i=1$,
impossible.  For $n$ even, $w$ is free on $i\le n/2$ and determined on the rest.  Now
$n=d-1$.
\end{proof}

\begin{proposition}[only the composite lifts]\label{prop:norev}
There is no invertible degree-preserving linear map $\mathcal J$ of the Fock space with
\[
   \mathcal J\,R_e^{W}\,\mathcal J^{-1}\;=\;R_e^{W\circ\rev}
   \qquad\text{for all } e\ge1 \text{ and all shape weights } W .
\]
The same holds with $\rev$ replaced by $\cpl$.
\end{proposition}

\begin{proof}
Suppose such a $\mathcal J$ exists.  Degree $0$ of the Fock space is spanned by
$s_\varnothing$, so $\mathcal J s_\varnothing=c\,s_\varnothing$ with $c\ne0$.

Take $e=3$ and let $\delta_{u}$ denote the weight supported on the single word $u$ with
value $1$.  Under \cite[Lem.~\texttt{lem:dict}]{Q149} the word $(1,0)$ is the ribbon with
row lengths $(2,1)$ read top first, and $(0,1)$ is the ribbon $(1,2)$; and
$\rev(1,0)=(0,1)$.  Hence the hypothesis gives
$\mathcal J\,R_3^{\delta_{(1,0)}}=R_3^{\delta_{(0,1)}}\,\mathcal J$.
Apply both sides to $s_\varnothing$.  A ribbon added to $\varnothing$ is a hook, so its
rows are weakly decreasing top to bottom; the shape $(1,2)$ is not, and the shape $(2,1)$
is.  Therefore $R_3^{\delta_{(0,1)}}s_\varnothing=0$ while
$R_3^{\delta_{(1,0)}}s_\varnothing=s_{(2,1)}$, so
$\mathcal J s_{(2,1)}=R_3^{\delta_{(0,1)}}(c\,s_\varnothing)=0$,
contradicting injectivity of $\mathcal J$.  For $\cpl$ the same two words work, since
$\cpl(1,0)=(0,1)$ as well.
\end{proof}

\begin{remark}[the shape of the answer]
$R_3^{\delta_{(0,1)}}$ is not the zero operator --- it sends $s_{(1)}\mapsto s_{(2,2)}$ ---
so the obstruction is not that one of the two tromino operators vanishes.  It is that they
vanish on \emph{different} vectors, and the vacuum separates them.  This is the smallest
live case: for $e\le2$ the involution $\rev$ is the identity on $\{0,1\}^{e-1}$, so nothing
can be seen there.
\end{remark}

\begin{corollary}[the honest accounting]\label{cor:accounting}
Of the $(\Z/2)^2$ available on words, exactly the subgroup $\langle T\rangle$ is implemented
by a Fock-space involution (Theorem~\ref{thm:A}, Proposition~\ref{prop:norev}); and what that
subgroup buys depends on where one stands on the locus \eqref{eq:locus}:
\begin{itemize}
\item on the nowhere-zero part it acts by the scalar $\kappa$, so
  $\operatorname{supp}(W)$ is already $T$-stable and no nonvanishing condition is paired with
  a different one: \emph{nothing is free};
\item on the single-shape part --- the part that is new, and that
  \cite[Lem.~8.2]{Q149} wrongly excluded --- it acts as the involution $w\mapsto w^{T}$, whose
  fixed points are the $2^{(d-1)/2}$ antipalindromes when $d$ is odd and none when $d$ is even.
  So the classification of single-shape commuting pairs need only be carried out on
  $T$-orbit representatives, and \emph{half of it is free}: the orbit counts are
  $1,2,2,6,6,18,24,60$ for $d=2,\dots,9$ against $0,2,4,8,12,28,48,104$ shapes.
\end{itemize}
In particular the mechanism of \cite[\texttt{prop:threefourths}]{T2020}
(Remark~\ref{rem:mechanism}) does have an analogue here --- but it is delivered by an
automorphism, not an antiautomorphism, and it delivers on exactly the half of the locus that
the superseded \cite[Thm.~\texttt{thm:main}]{Q149} did not see.
\end{corollary}

\begin{remark}[a terminology hazard, recorded]\label{rem:NBname}
\cite[Def.~\texttt{def:border}]{Q150c2} \emph{defines} ``non-overlapping'' to mean
\textnormal{(NB)}: every border length $p$ of $w$ has $w_{p+1}\ne w_{n-p}$.  This is weaker than
the standard meaning of ``non-overlapping'' (= unbordered = no proper nonempty border) in
combinatorics on words, and the two differ from $d=5$ on: at $d=5$ there are $8$ words
satisfying (NB) and only $6$ unbordered ones, the difference being $0101$ and $1010$.  The
counts $0,2,4,8,12,28,48,104$ of \cite[\S ver]{Q150c2} are the (NB) counts.  Any write-up that
merges these sources must not import the word; it must import the condition.
\end{remark}

\section{Verification}\label{sec:ver}

All checks are in \texttt{scratch/q151/} (\texttt{maya.py}, \texttt{sep.py}), exact
arithmetic over $\Q$ or $\Q(\gamma)$ via \textsc{SymPy}; no floating point.

\begin{description}
\item[V1 ($\Cop$ is transposition).]  For all $97$ partitions of size $\le9$,
  $\Cop(M(\lambda))=M(\lambda')$.  $0$ failures.

\item[V2 (Theorem~\ref{thm:A} at move level).]  For all $97$ partitions of size $\le9$ and
  all $e\le6$: the map $b\mapsto-1-b-e$ is a bijection from the legal rank-$e$ moves of
  $M(\lambda)$ to those of $\Cop M(\lambda)$; the occupancy word transforms by $u\mapsto
  u^{T}$; and the targets correspond under $\Cop$.  $2670$ moves, $0$ failures.

\item[V3 (the separator; Prop.~\ref{prop:variance}).]  With $W,\bar W$ \emph{independent
  generic symbolic} weights --- $2^{e-1}+2^{f-1}$ free symbols, deliberately off the
  commuting locus --- and for every partition $\lambda$ with $|\lambda|\le5$:
  \[
  \begin{array}{c|ccc}
    (e,f) & \Cop(R_eR_f)\Cop = R_e^\vee R_f^\vee & \Cop(R_eR_f)\Cop = R_f^\vee R_e^\vee
          & \#\{\lambda: \text{the two candidates differ}\}\\\hline
    (1,2) & 19/19 & 0/19 & 19/19\\
    (2,3) & 19/19 & 0/19 & 19/19\\
    (2,2) & 19/19 & 0/19 & 15/19\\
    (3,4) & 19/19 & 0/19 & 19/19\\
    (2,4) & 19/19 & 0/19 & 19/19
  \end{array}
  \]
  The last column is the positive control demanded by Remark~\ref{rem:sep}: the test is not
  constant in the direction it measures.

\item[V4 ($\kappa$; Thm.~\ref{thm:C}(ii)).]  For every $d\le7$, every $e\le12$ with $d\mid
  e$, and each choice $\gamma(d/2)=\pm1$ when $d$ is even, with $\gamma$ otherwise a free
  symbol on orbit representatives: $W^{\gamma}_e(u^{T})=\kappa\,W^{\gamma}_e(u)$ for all
  $u$, with $\kappa=(-1)^{e-1}\gamma(d/2)^{e/d}$.  $41$ triples $(d,e,\pm)$, $0$ failures.

\item[V5 (an untuned witness at $d=3$).]  $(e,f)=(3,6)$, $\gamma=(1,3,\tfrac13)$ on $\Z/3$
  --- a point of the commuting locus that is \emph{not} Murnaghan--Nakayama, and for which
  $\gamma\ne\gamma^{-1}$.  Then $[R_3^{W},R_6^{\bar W}]s_\lambda=0$ for all $12$ partitions
  $\lambda$ with $|\lambda|\le4$; likewise for $\gamma^{-1}$.  Directly:
  $W\circ\rev=W^{\gamma^{-1}}$ and $\bar W\circ\rev=\bar W^{\gamma^{-1}}$, while
  $W\circ T=+W$ and $\bar W\circ T=-\bar W$, matching $\kappa=(-1)^{2}=1$,
  $\bar\kappa=(-1)^{5}=-1$.  This is the separator of Theorem~\ref{thm:C}(iii): reversal
  moves the point, the composite does not.

\item[V6 (Prop.~\ref{prop:norev}).]  $R_3^{\delta_{(0,1)}}s_\varnothing=0$,
  $R_3^{\delta_{(1,0)}}s_\varnothing=s_{(2,1)}$, $R_3^{\delta_{(0,1)}}s_{(1)}=s_{(2,2)}$.

\item[V7 (Lemma~\ref{lem:otimes}).]  With $\tau$ a generic symbolic rank-$d$ weight: for all
  $d\le4$, $k\le3$ with $kd-1\le9$, and $g\in\{\rev,\cpl,T\}$,
  $(\tau^{\otimes k})(g\,y)=(\tau\circ g)^{\otimes k}(y)$ for every word $y$.  $1368$
  identities, $0$ failures.

\item[V8 (Lemma~\ref{lem:2Bstable}).]  Exhaustive over $\mathbb F_5$: all $25$ solutions of
  $(2\mathrm B^{*}_3)$ and all $81$ of $(2\mathrm B^{*}_4)$ remain solutions after $\rev$,
  $\cpl$ and $T$.  $0$ failures.

\item[V9 (Thm.~\ref{thm:D}(b); the separator for ``$T$ acts by a scalar'').]  Exhaustive over
  $\mathbb F_5$ and $\mathbb F_7$, $m=3,4$: every nonzero $(2\mathrm B^{*}_3)$ weight satisfies
  $\sigma\circ T\in F^{\times}\sigma$, but at $m=4$ exactly $32$ of $80$ (over $\mathbb F_5$)
  and $72$ of $168$ (over $\mathbb F_7$) do not --- the failures being the single-shape weights
  on non-$T$-symmetric words, e.g.\ $\sigma=\mathbb 1_{110}$ with
  $\sigma\circ T=\mathbb 1_{100}$.  \textbf{This is the check that refuted the first version of
  this section}, which had asserted the scalar action on the whole locus on the strength of
  the superseded \cite[Thm.~\texttt{thm:main}]{Q149}.

\item[V10 (the (NB) family and its $T$-orbits).]  Direct enumeration for $d\le9$:
  $\{w:\mathbb 1_w$ satisfies $(2\mathrm B^{*}_d)\}$ agrees with (NB) of
  \cite[Def.~\texttt{def:border}]{Q150c2} in all $\sum_d 2^{d-1}$ cases, with counts
  $0,2,4,8,12,28,48,104$ for $d=2,\dots,9$ --- matching \cite[\S ver]{Q150c2}.  The family is
  $\rev$-, $\cpl$- and $T$-stable; $\mathrm{Fix}(T)$ inside it equals the antipalindromes
  exactly, with counts $0,2,0,4,0,8,0,16$; the $T$-orbit counts are $1,2,2,6,6,18,24,60$.
  Separately, the (NB) family strictly contains the \emph{unbordered} words from $d=5$ on
  ($8$ versus $6$), which is Remark~\ref{rem:NBname}.
\end{description}

\section{What is not proved}

\begin{enumerate}
\item Proposition~\ref{prop:norev} rules out invertible \emph{degree-preserving} linear
implementations of $\rev$.  It does not rule out an implementation by a map that mixes
degrees, nor one defined only on a subquotient.  I believe no such map exists but have not
proved it; the hypothesis is used exactly once, to force $\mathcal J
s_\varnothing\in\F^\times s_\varnothing$.

\item Theorem~\ref{thm:D} treats two parts of the locus \eqref{eq:locus}: the nowhere-zero
seeds and the single-shape seeds.  These do not exhaust it --- a seed $\tau$ may have zeros
without being supported on one word --- and for the intermediate seeds I have not determined
when $\tau\circ T\in F^{\times}\tau$.  What Check~V9 does establish is that the scalar action
fails somewhere, which is all Corollary~\ref{cor:accounting} needs; the precise
$T$-fixed locus of \eqref{eq:locus} is open.

\item Corollary~\ref{cor:accounting}'s ``half of it is free'' is an accounting statement
about a classification, not a proof that any particular proof shortens.  Concretely: for each
$T$-orbit $\{w,w^{T}\}$ of size $2$, Theorem~\ref{thm:A} gives the second commuting pair from
the first by conjugation; it does not give the \emph{verification} that $w$ satisfies (NB) for
free, since that was never the expensive half.

\item The Leonard-system analogy is now closed in the negative for $\dagger$, but not for
the rest of \cite{T2020}: the finding of \texttt{sources.json} that PA2 (nonvanishing) and
PA5 (rigidity) appear there as \emph{independent} conditions remains the live warning, and
is untouched by anything above.
\end{enumerate}

\begin{thebibliography}{9}
\bibitem{Q149} Clio, \emph{Shape weights on ribbons: the commuting locus is the
transpose-equivariant $\gcd(e,f)$-periodic characters},
\texttt{proofs/2026-09-15-c1-Q149-shape-weights.tex}, 15 September 2026.
\bibitem{Q150c1} Clio, \emph{Q150: the nonvanishing lemma is false},
\texttt{proofs/2026-09-16-c1-Q150-nonvanishing.tex}, 16 September 2026.
\bibitem{Q150c2} Clio, \emph{Q150 (H2): single-shape commuting operators},
\texttt{proofs/2026-09-16-c2-Q150-H2-single-shape.tex}, 16 September 2026.
\bibitem{T2020} P.~Terwilliger, \emph{Notes on the Leonard system classification},
\texttt{arXiv:2003.09668}; Graphs and Combinatorics \textbf{37} (2021) 1687--1748.
Definitions and statements above are quoted from the arXiv LaTeX source
(\texttt{arxiv.org/e-print/2003.09668}), \S4 (\texttt{lem:antiaut}) and
\texttt{prop:threefourths}.
\end{thebibliography}

\end{document}
